% !TeX root = ../教材题库_全部条目_转写工作稿.tex
% Source: XJ-B1-C05-EX-P182-Y0945
% 第5章 / 5.3 三角函数的图象与性质 / 教材 p.182 / PDF 第189页
% OCR未校对：公式与文字均可能有误；未完成数学TeX转写；禁止直接发给学生
% 原题及相邻原解截图（核对用）：../题目原图/XJ-B1-C05-EX-P182-Y0945-01.webp
% 本文件只在转写完成后作为题干使用；OCR 条目尚不是可靠题干。
\begin{bankitem}{5.3 例题 例题2}{第5章 / 5.3 三角函数的图象与性质 / 教材 p.182 / PDF 第189页}{XJ-B1-C05-EX-P182-Y0945}{OCR未校对：公式与文字均可能有误；未完成数学TeX转写；禁止直接发给学生}
\begin{ocrdraft}
[教材p.182 / PDF第189页]
\par
00求下列函数的最大值及取得最大值时自变量的集合：
\par
（D3，=2sin2，x\ensuremath{\text{[OCR:EUR]}}R；
\par
（2）3，一2一cos，xGR.
\par
解（1）令之一2使函数y=2sin之，之eR取得最大值的之的集合是
\par
之之一一十2差丌，k\ensuremath{\text{[OCR:EUR]}}Z
\par
由2一之一一十2差丌，得一一十差丌，k\ensuremath{\text{[OCR:EUR]}}Z．
\par
因此，使函数3，一2sin2厶xeR取得最大值的工的集合是x一一十差丌，kez
\par
最大值是2．
\par
（2）当函数y一2一cos“eR取得最大值时，cos取
\par
最小值一1，此时工的集合是
\par
\{工一丌十2差丌，k\ensuremath{\text{[OCR:EUR]}}Z\}
\par
你能求出例2中各函
\par
因此，使函数3-2一cos工，\ensuremath{\text{[OCR:EUR]}}R取得最大值的x的集合
\par
数取得最小值时自变量工
\par
是0应一丌十2差丌，k\ensuremath{\text{[OCR:EUR]}}Z\}，最大值是3．
\par
的集合吗？
\end{ocrdraft}
\end{bankitem}
